\section{Wprowadzenie: pytanie menedżerskie}

Decyzja inwestycyjna zarządu polskiej firmy w roku 2026 nie zapada w podręcznikowym stanie ogólnej równowagi. Zapada w środowisku, w którym wynagrodzenia, koszt finansowania, kurs PLN/EUR, ceny importowanych komponentów, dostępność kredytu i oczekiwania inflacyjne zmieniają się równocześnie. Dla zarządu i dyrektora finansowego pytanie o Bezwarunkowy Dochód Podstawowy nie jest więc wyłącznie pytaniem o redystrybucję. Jest pytaniem o marżę, płynność, inwestycje, strukturę kosztów i zdolność obsługi długu w kolejnych kwartałach po wprowadzeniu takiego programu.

W debacie publicznej BDP najczęściej występuje jako odpowiedź na automatyzację: skoro sztuczna inteligencja, robotyzacja i oprogramowanie wypierają część pracy, państwo powinno rozważyć dochód podstawowy jako zabezpieczenie społeczne \parencite{autor2015,brynjolfsson2014,acemoglu2018,acemoglu2020,graetzMichaels2018,freyOsborne2017,korinekStiglitz2019,susskind2020}. Taki kierunek rozumowania jest intuicyjny, ale jednostronny. Ta praca wychodzi od odwrotnej możliwości: BDP może nie tylko amortyzować skutki automatyzacji, lecz także zmieniać bodźce, które prowadzą firmy do automatyzacji. Innymi słowy, transfer publiczny może stać się jednym z mechanizmów przyspieszających substytucję pracy kapitałem.

Mechanizm ten nie jest jednak prosty. BDP może podnosić płacę rezerwową, czyli minimalny poziom wynagrodzenia, przy którym część gospodarstw domowych akceptuje zatrudnienie. To zwiększa presję kosztową po stronie firm. Równocześnie BDP zwiększa dochód gospodarstw domowych i może wspierać popyt, co poprawia warunki sprzedażowe przedsiębiorstw. Ten sam transfer powiększa jednak wydatki publiczne, deficyt i dług, a przez inflację, reakcję banku centralnego, rentowności obligacji, kurs walutowy i koszt długu korporacyjnego wpływa na koszt kapitału. Firma znajduje się więc między dwoma ruchami cen względnych: pracą, która może stać się droższa, oraz kapitałem, który także może stać się droższy.

Z tego powodu główne pytanie pracy nie brzmi: czy BDP jest dobry albo zły. Pytanie jest bardziej praktyczne: w jakim zakresie BDP może działać jako katalizator inwestycji w AI i automatyzację, a od którego momentu interakcja presji płacowo-rezerwowej oraz kanału fiskalno-finansowego zaczyna ograniczać zdolność firm do sfinansowania tej reakcji. Łączy ono makroekonomię z zarządzaniem, bo przekłada się na decyzje o długu, nakładach inwestycyjnych, technologii, zatrudnieniu i zabezpieczeniu ekspozycji walutowej.

Punktem odniesienia jest raport Międzynarodowego Funduszu Walutowego \textit{Fiscal Monitor: Tackling Inequality} z 2017 r. Raport MFW jest w tej pracy traktowany jako punkt odniesienia dla analizy redystrybucji, efektywności i dobrobytu w określonej ramie modelowej. W przypadku Polski wykorzystuje analizę statyczną opartą na danych LIS z 2013 r. To ujęcie pomaga zakotwiczyć skalę problemu, ale nie wyczerpuje problemu menedżerskiego. Dla firmy ważne jest nie tylko to, gdzie gospodarka może znaleźć się w nowej równowadze, lecz także to, co dzieje się po drodze: z płynnością, bilansem, kredytem, kursem walutowym, bankructwami, inwestycjami i heterogenicznymi reakcjami przedsiębiorstw.

Dlatego w pracy wykorzystano \amor{}, autorski wykonywalny model stock-flow consistent agent-based gospodarki polskiej rozwijany w ramach \bbg{}. Model jest skalibrowany do stanu Polski na 2026-04-30, a eksperyment BDP traktowany jest jako kontrfaktyczna ścieżka od tego punktu startowego. \amor{} nie opiera się na reprezentatywnym agencie: operuje na heterogenicznych gospodarstwach domowych, firmach, bankach, sektorze publicznym, banku centralnym i sektorze zewnętrznym. Każdy przepływ pieniężny w modelu musi mieć drugą stronę księgową. Hipoteza zostaje więc zapisana w mechanizmie, uruchomiona i porównana ze ścieżką bez BDP.

Wynik eksperymentu doprecyzowuje początkową intuicję. Mocna wersja hipotezy, zgodnie z którą wyższy BDP powinien prowadzić do coraz większej automatyzacji, nie znajduje potwierdzenia. W wariancie centralnym adopcja AI/hybrid rośnie do okolic scenariusza BDP 2000~PLN, po czym spada przy wyższych poziomach transferu. Oznacza to, że BDP może być katalizatorem automatyzacji tylko warunkowo: dopóki presja kosztowa i popytowa nie ograniczają zdolności sektora prywatnego do finansowania inwestycji. Dla zarządu próg absorpcji jest więc zmienną strategiczną: automatyzacja nie jest tylko wyborem technologicznym, lecz decyzją finansową podejmowaną pod ograniczeniem płynności, zadłużenia i ryzyka makroekonomicznego.
